% ==================== 2 实验部分 ====================
\section{实验部分}

\subsection{样品与实验变量}

本文使用的 8 组样品均为 TiO$_2$ 纳米粉体，样品编号 200、300、$\ldots$、900 对应不同的热处理温度（$^\circ$C）。样品来自同一批前驱体，热处理后保存在干燥器中；除温度外，其余制备条件尽量保持一致，以便把衍射谱的差异归因于热处理温度这一实验变量。

保温时间、升降温速率和气氛都按实验记录执行，样品编号只保留温度信息。因此后文说“200--600 $^\circ$C 以锐钛矿为主、700 $^\circ$C 出现金红石”时，指的是样品编号对应的热处理温度，不是严格的热力学转变温度。

\subsection{仪器与测试参数}

测试在 D-MAX 3 型 X 射线衍射仪上完成，使用 Cu 靶 K$\alpha$ 辐射（$\lambda=1.54056$ \AA），管压 40 kV、管流 30 mA。扫描范围 $2\theta=20^\circ$--$80^\circ$，步宽 0.02$^\circ$，每组谱 3000 个数据点。8 组样品在同一台仪器、同一组狭缝和同一扫描参数下连续测试，以减少仪器状态变化带来的系统误差。主要参数列于表~\ref{tab:params}。

衍射仪由 X 射线光管、测角仪和探测器三部分组成。光管中的灯丝发射电子，经高压加速后轰击 Cu 靶，产生 Cu K$\alpha$ 射线；射线经发散狭缝和 Soller 狭缝准直后照到样品上。测角仪让样品和探测器按 $\theta$--$2\theta$ 关系联动，探测器把不同 $2\theta$ 方向的衍射光子转换成计数。光管的能量转换效率不高，大部分功率以热的形式耗散，所以阳极靶必须通循环冷却水；这也是开机前要先开冷却水的原因。

\begin{table}[H]
  \centering
  \caption{X 射线衍射测试参数}
  \label{tab:params}
  \begin{tabular}{@{}ll@{}}
    \toprule
    项目 & 参数 \\
    \midrule
    仪器 & D-MAX 3 型 X 射线衍射仪 \\
    靶材与波长 & Cu K$\alpha$，$\lambda=1.54056$ \AA \\
    管压 / 管流 & 40 kV / 30 mA \\
    扫描范围 & $2\theta=20^\circ$--$80^\circ$ \\
    步宽 / 数据点数 & 0.02$^\circ$ / 3000 \\
    样品 & TiO$_2$，200--900 $^\circ$C 热处理，共 8 组 \\
    \bottomrule
  \end{tabular}
\end{table}

\subsection{样品制备与测试流程}

制样对数据质量的影响很大。样品在玛瑙研钵中磨到用手指捻时没有颗粒感，填入样品槽后用毛玻璃片压实，让上表面与槽沿齐平；装好后擦净槽外和台面上的散落粉末，避免污染测角台。具体手法和常见失误在第 5 节汇总。

开机先开循环冷却水，等水温降到约 19 $^\circ$C 再打开衍射仪钥匙和 XD-3 软件。管压、管流先设在 15 kV、6 mA 预热几分钟，再缓慢升到工作值，一步到位容易损伤灯丝和靶。样品台装到测角台后，用软件中的“校读”功能校准测角仪初始位置，确认样品表面没有碰到光路，关好防护门再开始扫描。本次设定起始角 20$^\circ$、终止角 80$^\circ$、步宽 0.02$^\circ$，扫描结束后立即存为 txt 文件；换样时重复“装样—校读—关门—扫描”这几步。

关机按相反顺序：先把管压、管流降到 15 kV、6 mA，等约 5 min 再关 X 射线；约 10 min 后关循环水，最后左旋钥匙关闭衍射仪，并填写仪器使用记录。详细步骤见实验讲义\cite{shi2024xrdlab}。
